\section*{अध्याय \ref{ch:b2:himalaya-andes} --- टक्कर-पट्टियाँ और एंडीज़-प्रकार की पट्टियाँ: एक तुलना}

\begin{solution}{exo:b2:himalaya-andes:1}
(क) $r = h\rho_c/(\rho_m - \rho_c) = 3.0\times2700/600 = \qty{13.5}{km}$। (ख) $T = 35 + 3.0 +
13.5 = \qty{51.5}{km}$, या $35 + 5.5\times3.0$। (ग) $h = \Delta T(\rho_m - \rho_c)/\rho_m =
10\times600/3300 = \qty{1.8}{km}$।
\end{solution}

\begin{solution}{exo:b2:himalaya-andes:2}
(क) सिंधु-गंगा अग्रगर्त; मुख्य अग्र अधिक्षेप; उप-हिमालय (शिवालिक); मुख्य सीमा
अधिक्षेप; लघु हिमालय; मुख्य केंद्रीय अधिक्षेप; बृहत् हिमालयी अनुक्रम; दक्षिण तिब्बत
विलगन तंत्र; टेथिस हिमालयी अनुक्रम; फिर सिंधु--त्सांगपो सीवन और एशियाई
भूपर्पटी। (ख) दक्षिण तिब्बत विलगन तंत्र। (ग) श्रेणी उत्तर की ओर बढ़ती है: अकायांतरित
\omterm{def:b2:alpine-orogen:foreland}{मोलास}, मध्यम-श्रेणी का लघु हिमालय, उच्च-श्रेणी और आंशिक रूप से गलित बृहत् हिमालयी
अनुक्रम; STDS के ऊपर यह अचानक गिरकर अकायांतरित या दुर्बल रूप से कायांतरित
टेथिस अवसादों तक आ जाती है, क्योंकि सामान्य भ्रंश ने बीच की चट्टानों को हटा दिया है।
\end{solution}

\begin{solution}{exo:b2:himalaya-andes:3}
(क) $196.5\ \text{m}/9200\ \text{yr} = \qty{0.02136}{m}$ प्रति वर्ष, यानी लगभग \qty{21.4}{mm} प्रति वर्ष।
(ख) यह \qty{20}{\percent} अधिक है, GPS के \qty{17.8}{mm} प्रति वर्ष से, पर उसी परिमाण का:
एक दर नौ हज़ार वर्षों पर औसत है, दूसरी बीस वर्षों पर।
(ग) पर्वतमाला के आर-पार लगभग सारा लघुकरण अग्र पर खपता है: MBT, MCT और
MFT के उत्तर के अन्य भ्रंश, इस देशांतर पर, व्यवहारतः निष्क्रिय हैं; वे सभी
उसी MHT में जड़ जमाते हैं, जो MFT पर ऊपर आता है।
\end{solution}

\begin{solution}{exo:b2:himalaya-andes:4}
(क) टक्कर: सीवन लुप्त महासागर को चिह्नित करता है। (ख) एंडीज़-प्रकार: चाप को नीचे जाता हुआ
स्लैब चाहिए। (ग) सामान्यतः टक्कर: मोटी हुई, गर्म हुई भूपर्पटी के भूपर्पटीय गलित (ये
मोटी हुई एंडीज़ भूपर्पटी में भी मिलते हैं, पर वहाँ चाप के \omterm{def:b2:magma-emplacement:magma}{मैग्मा} प्रबल हैं)। (घ) एंडीज़-प्रकार।
(ङ) दोनों: शिवालिक और उप-एंडीज़ अग्रभूमि। (च) दोनों: हिमालय में MHT,
एंडीज़ में क्षेपण-अंतरापृष्ठ।
\end{solution}

\begin{solution}{exo:b2:himalaya-andes:5}
(क) पूर्व में $17.8/45 = \qty{40}{\percent}$, पश्चिम में $20.5/45 = \qty{46}{\percent}$:
लगभग आधा। (ख) तिब्बत और बाक़ी एशिया के स्थूलन और भ्रंशों में, और
बड़े स्ट्राइक-स्लिप भ्रंशों के अनुदिश खंडों के पूर्व की ओर बहिर्वेधन में। (ग) $45 - 17.8 =
\qty{27.2}{mm}$ प्रति वर्ष, यानी \qty{272}{km}, \qty{10}{Myr} में; $45 - 20.5 = \qty{24.5}{mm}$ प्रति
वर्ष, यानी \qty{245}{km}।
\end{solution}

\begin{solution}{exo:b2:himalaya-andes:6}
(क) $r = 4.5\times3.7 = \qty{16.65}{km}$; $\Delta F = 2700\times9.8\times3700\times(35\,000 +
1850 + 8325) = \qty{4.42e12}{N.m^{-1}}$। (ख) तिब्बत: \qty{6.12e12}{N.m^{-1}}; आल्टिप्लानो
के पास उसका \qty{72}{\percent} है। (ग) $r = 4.5h$ के साथ, $\Delta F = \rho_c g\,(Hh + 2.75h^2)$:
$h \ll H$ के लिए पद $Hh$ प्रबल है और $\Delta F$ $h$ के अनुपाती है; ऊँचे पठार के लिए
$h^2$ वाला पद तेज़ी से बढ़ता है। ऊँचाई के हर अगले किलोमीटर के लिए किनारों पर उत्तरोत्तर अधिक धक्का चाहिए,
जब तक उपलब्ध बल पठार को और नहीं उठा पाते; तब वह ऊपर उठने के बजाय
फैलता है।
\end{solution}

\begin{solution}{exo:b2:himalaya-andes:7}
(क) $2700\times9.8\times8848.86/2 = \qty{1.17e8}{Pa} = \qty{117}{MPa}$। (ख)
\qty{100}{MPa} पर: $0.85\times100 = \qty{85}{MPa}$; \qty{200}{MPa} पर: $0.85\times200 = 50 +
0.6\times200 = \qty{170}{MPa}$, दोनों शाखाएँ सहमत हैं। (ग) $h = 2\sigma_d/(\rho g) =
2\times1.70\times10^8/(2700\times9.8) = \qty{12.8}{km}$: परिमाण की एक कोटि, सबसे ऊँचे शिखरों की
कोटि की, उनसे बहुत ऊपर नहीं।
\end{solution}

\begin{solution}{exo:b2:himalaya-andes:8}
(क) $11\times111.19\times\cos\ang{30} = \qty{1059}{km}$। (ख) $9\times10^{-3}/1.059\times10^{6}
= \num{8.5e-9}$ प्रति वर्ष (\num{6.6e-9} और \num{10.4e-9} के बीच, 7 और \qty{11}{mm} प्रति वर्ष के लिए),
यानी $8.5\times10^{-9}/3.156\times10^{7} = \qty{2.7e-16}{s^{-1}}$। (ग) दोनों परास
अतिव्याप्त हैं: $(12.4 \pm 4.7) \times 10^{-9}$ 7.7 से \num{17.1e-9} तक फैला है। नेपाल के उत्तर का दक्षिणी पठार
पूरी पट्टी जितनी तेज़ी से, या उससे थोड़ा तेज़, खिंचता है।
\end{solution}

\begin{solution}{exo:b2:himalaya-andes:9}
(क) $0.9/6.6 = \qty{14}{\percent}$। (ख) $M_0 = 10^{1.5\times8.5 + 9.1} = 10^{21.85} =
\qty{7.1e21}{N.m}$। (ग) $7.1\times10^{21}/6.6\times10^{19} = 107$: नेपाल के नीचे ऐसा एक भूकंप लगभग
हर 110 वर्ष में, यानी इतिहास जितना अंकित करता है उससे कहीं अधिक बार; या तो कहीं बड़े भूकंप
आने वाले हैं, या घाटे का एक भाग अन्य ढंग से मुक्त होता है।
\end{solution}

\begin{solution}{exo:b2:himalaya-andes:10}
(क) प्रतिपूर्ति-गहराई आल्टिप्लानो भूपर्पटी के आधार पर लीजिए, समुद्र-तल से $61.5 - 3.7 =
\qty{57.8}{km}$ नीचे। पूना भूपर्पटी $50.5 - 4.4 = \qty{46.1}{km}$ तक पहुँचती है, इसलिए
पूना-स्तंभ में उस गहराई के ऊपर \qty{11.7}{km} प्रावार है: $2800\times61.5 =
\rho_P\times50.5 + 3300\times11.7$। (ख) $\rho_P = (172\,200 - 38\,610)/50.5 =
\qty{2645}{kg.m^{-3}}$। (ग) अंतर, \qty{155}{kg.m^{-3}}, उद्धृत 150 से \qty{200}{kg.m^{-3}} के
निचले सिरे पर है: केवल भूपर्पटी के लिए असामान्य रूप से हल्की पूना भूपर्पटी चाहिए होगी।
उसके नीचे हल्का स्तंभ यह काम अधिक सरलता से करता है: विस्तरीकरण से हटा प्रावार-स्थलमंडल,
जिसका स्थान गर्म, हल्के दुर्बलमंडल ने ले लिया।
\end{solution}

\begin{solution}{exo:b2:himalaya-andes:11}
(क) $g = 105/\tan\ang{30} = \qty{181.9}{km}$। (ख) $g = 90/\tan\ang{30} + 15/\tan\ang{4} =
155.9 + 214.5 = \qty{370.4}{km}$: चाप \qty{188.5}{km} भीतर की ओर खिसकता है। (ग) स्लैब प्लेट के नीचे कभी
\qty{105}{km} तक नहीं पहुँचता, उन \qty{300}{km} में, और उसके ऊपर कोई गर्म प्रावार-कील
नहीं होती: कोई \omterm{def:b2:magma-emplacement:magma}{मैग्मा} नहीं बनता और चाप तब तक मृत रहता है जब तक स्लैब फिर से, और पूर्व में, नहीं धँसता।
\end{solution}

\begin{solution}{exo:b2:himalaya-andes:12}
(क) \qty{21}{mm} प्रति वर्ष का अर्थ है प्रति मिलियन वर्ष \qty{21}{km}: $21\times37.5 =
\qty{787.5}{km^3}$ पर्वतमाला के प्रति किलोमीटर, प्रति मिलियन वर्ष। (ख) $300/787.5 =
\qty{38}{\percent}$। (ग) पर्वतमाला और दक्षिणी तिब्बत के स्थूलन में, भारत के साथ तिब्बत के नीचे
ले जाई गई भूपर्पटी में, और उस पदार्थ में जिसे MHT उत्तर की ओर आगे बढ़ाता है।
\end{solution}

\begin{solution}{pb:b2:himalaya-andes:1}
\textbf{1.} \ang{28.0} से \ang{37.0}~N तक सात नमूने: $34\,550/7 = \qty{4936}{m}$, $h_T
= \qty{4.9}{km}$।
\textbf{2.} \ang{-69} से \ang{-65} तक पाँच नमूने: $18\,650/5 = \qty{3730}{m}$, $h_A =
\qty{3.7}{km}$, जो आल्टिप्लानो के लिए उद्धृत मान है; पूना, और दक्षिण में, अधिक ऊँचा है।
\textbf{3.} तिब्बत: कम से कम $9\times111.19 = \qty{1001}{km}$; आल्टिप्लानो: कम से कम $4\times
104.5 = \qty{418}{km}$।
\textbf{4.} $(5000 - 70)/(1.5\times111\,190) = 0.0296$: \qty{3.0}{\percent}।
\textbf{5.} एंडीज़ का प्रोफ़ाइल पठार के पश्चिम में महासागर-तल और खाई में, समुद्र-तल से
किलोमीटरों नीचे, आरंभ होता है; तिब्बत वाला समुद्र-तल के निकट एक महाद्वीपीय मैदान पर आरंभ होता है,
और ऊँची तारिम द्रोणी पर समाप्त होता है।
\textbf{6.} $r = 4.5\times4.9 = \qty{22.05}{km}$; $T_T = 35 + 4.9 + 22.05 = \qty{62.0}{km}$।
\textbf{7.} $r = 4.5\times3.7 = \qty{16.65}{km}$; $T_A = 35 + 3.7 + 16.65 = \qty{55.4}{km}$।
\textbf{8.} तिब्बत: \qty{62.0}{km}, जबकि मापा गया मान 70 से \qty{75}{km} है, यानी \qty{10.5}{km} कम, मध्य
मान \qty{72.5}{km} से (\qty{14}{\percent}); आल्टिप्लानो: \qty{55.4}{km}, जबकि मापा गया मान 59 से
\qty{64}{km} है, यानी \qty{6.1}{km} कम, \qty{61.5}{km} से (\qty{10}{\percent})। तिब्बत अधिक दूर है।
\textbf{9.} $72.5 - 35 = 4.9\times3300/(3300 - \rho_c)$: $3300 - \rho_c = 16\,170/37.5 =
431.2$, $\rho_c = \qty{2869}{kg.m^{-3}}$।
\textbf{10.} $26.5 = 3.7\times3300/(3300 - \rho_c)$: $3300 - \rho_c = 460.8$, $\rho_c =
\qty{2839}{kg.m^{-3}}$।
\textbf{11.} भार का एक भाग प्लेट के झुकने से वहन होता है, किसी स्थानीय जड़ से नहीं;
सामान्य भूपर्पटी शायद \qty{35}{km} मोटी न रही हो; नीचे का प्रावार मान्य से हल्का या सघन
हो सकता है; भूकंपीय मोहो शायद घनत्व के अर्थ में भूपर्पटी का आधार न हो।
\textbf{12.} तिब्बत की भूपर्पटी का आधार $72.5 - 4.9 = \qty{67.6}{km}$ पर है, भारत का
\qty{37.5}{km} पर: $t = \qty{30.1}{km}$। $\Delta g = 2\pi\times6.674\times10^{-11}\times370\times
30\,100 = \qty{4.67e-3}{m.s^{-2}}$, यानी भारत के सापेक्ष \qty{-467}{mGal}, जो
प्रेक्षित \qty{460}{mGal} के निकट है।
\textbf{13.} $W_0 = 440\times60/35 = \qty{754.3}{km}$; $S = \qty{314.3}{km}$।
\textbf{14.} $(314.3 - 250)/314.3 = \qty{20}{\percent}$, वही कमी जो मध्य
एंडीज़ में पाई गई।
\textbf{15.} $0.55\times314.3 = \qty{173}{km}$, जो संतुलित काटों के \qty{250}{km} से कम है,
और तब वे पर्याप्त से अधिक होंगी।
\textbf{16.} \omterm{def:b2:alpine-orogen:decollement}{डेकोलमां} के ऊपर छोटे हुए आवरण का मेल उसके नीचे, और पश्चिम में, आधारशैल के लघुकरण से
होना चाहिए, क्योंकि पट्टी के आर-पार दोनों को एक ही लंबाई खोनी है; आधारशैल का
लघुकरण पूर्वी कॉर्डिलेरा और पठार के नीचे भूपर्पटी को मोटा करता है।
\textbf{17.} $314.3/30 = \qty{10.5}{mm}$ प्रति वर्ष, यानी $10.5/72 = \qty{15}{\percent}$
अभिसरण का।
\textbf{18.} यह क्षेपण में खपता है: नाज़्का प्लेट महाद्वीप के नीचे धँसती है।
\textbf{19.} $r = \qty{22.05}{km}$; $\Delta F = 2700\times9.8\times4900\times(35\,000 + 2450
+ 11\,025) = \qty{6.28e12}{N.m^{-1}}$।
\textbf{20.} $2700\times9.8\times3700\times45\,175 = \qty{4.42e12}{N.m^{-1}}$; अनुपात
0.70 है।
\textbf{21.} $6.28\times10^{12}/10^{5} = \qty{6.3e7}{Pa}$, यानी लगभग \qty{63}{MPa}।
\textbf{22.} $9\times10^{-3}/1.06\times10^{6} = \num{8.5e-9}$ प्रति वर्ष।
\textbf{23.} भारत का धक्का उत्तर--दक्षिण कार्य करता है और ऊँचे पठार में
पठार के अपने भार से संतुलित होता है: उत्तर--दक्षिण पठार अब और मोटा नहीं हो सकता, इसलिए ऊर्ध्वाधर प्रतिबल,
यानी भार, महत्तम है; लघुत्तम प्रतिबल पूर्व--पश्चिम है, एशिया के मुक्त किनारों की ओर,
और पठार उसी ओर, उत्तर--दक्षिण दरारों पर, फैलता है।
\textbf{24.} \textbf{\qty[text-series-to-math=true]{62}{km}}, यानी लगभग \qty{86}{\percent}, मापे गए
\qty{72.5}{km} का।
\end{solution}
